\documentclass[12pt,letterpaper]{article}

% ---- Packages ----
\usepackage[utf8]{inputenc}
\usepackage[T1]{fontenc}
\usepackage{amsmath,amssymb}
\usepackage{newtxtext,newtxmath}  % Times New Roman text + math (replaces mathptmx)
\usepackage[margin=1in]{geometry}
\usepackage{setspace}
\usepackage{titlesec}
\usepackage{enumitem}
\usepackage{xr-hyper}
\usepackage{hyperref}
\externaldocument{1 - General Overview}
\externaldocument{3 - Gravity}
\externaldocument{5 - Strong and Weak Forces}
\externaldocument{7 - Quantum Mechanics}
\usepackage{microtype}
\usepackage{booktabs}

% ---- Section formatting ----
\titleformat{\section}{\large\bfseries\sffamily}{\thesection.}{0.5em}{}
\titleformat{\subsection}{\normalsize\bfseries\sffamily}{\thesubsection}{0.5em}{}
\titleformat{\subsubsection}{\normalsize\bfseries}{\thesubsubsection}{0.5em}{}

% ---- Spacing ----
\setstretch{1.15}
\setlength{\parskip}{0.5em}
\setlength{\parindent}{0em}

% ---- Custom commands ----
\newcommand{\rhoa}{\rho_{\!A}}
\newcommand{\Ka}{K_{\!A}}
\newcommand{\vin}{v_{\mathrm{in}}}
\newcommand{\rs}{r_s}
% \hbar is already defined in LaTeX; no redefinition needed

\begin{document}

\begin{center}
{\LARGE\bfseries Unifying Theory of Dynamic Aether Mechanics:\\
Cosmology}\\[1em]
{\large Erik Otto}
\end{center}

\vspace{1em}

\begin{abstract}
Within the Dynamic Aether Mechanics framework, cosmological phenomena emerge
from the three independently varying properties of the Aether medium (density,
pressure, temperature).  This paper explains cosmological redshift as real Doppler
shift from thermal expansion of cosmic voids, dark energy as the thermal pressure
of hot void Aether, and dark matter as the two-fluid first-sound correction
(the outward increase in $c_1$ from the varying normal fraction provides
additional centripetal acceleration through the acoustic metric).  Black holes are identified as sonic horizons where Aether inflow
velocity reaches $c$, providing a concrete physical mechanism for Hawking radiation
via asymmetric transient binding at the horizon boundary.  Neutron stars serve as
natural laboratories for extreme-density Aether physics, with pulsar glitches
arising from quantized vortex avalanches.
\end{abstract}

\vspace{1em}

\bigskip
\noindent\textit{This paper is part of a series presenting the Unifying Theory of
Dynamic Aether Mechanics.  For the foundational concepts of Aether binding,
the three independent properties (density, pressure, temperature), and the
unification of force constants, see Paper~1: General Overview \cite{otto-overview}.}
\bigskip

\section{Cosmological Effects}
% ============================================================

This is where the interplay between Aether's independently varying
properties---particularly temperature and the two-fluid normal fraction---determines
the large-scale structure and behavior of the universe.

Over cosmological timescales, two processes have shaped the thermal state of the Aether:

\begin{enumerate}[leftmargin=2em]
  \item Permanent binding (driven by the Dynamic Casimir Effect in high-energy environments)
    has created matter, locking thermal energy into mass and cooling the Aether in
    matter-containing regions.

  \item Stars and other luminous matter radiate energy outward into surrounding voids, where
    no mass formation absorbs it.\ This accumulated radiation, together with viscous
    dissipation from the Hubble flow (Section~\ref{sec:dark-energy}), heats the Aether in
    voids far above the temperature near matter.
\end{enumerate}

The Aether density $\rhoa$ is essentially uniform ($\delta\rhoa/\rhoa \sim
10^{-4}$; see Section~\ref{sec:hubble-tension}), because $\Ka \sim
10^{28}$~J/m$^3$ resists compression far more effectively than any cosmological
force can drive it.\ The cosmologically relevant variable is \emph{temperature}:
matter-containing regions have cold, low-normal-fraction Aether, while deep space
voids have hot, high-normal-fraction Aether.\ These two thermal environments
produce different large-scale effects.


\subsection{Cosmological Redshift}

The high temperature and thermal pressure of Aether in voids produces a real, physical
outward force on the matter at void boundaries.\ Since matter regions have low thermal
pressure (cold, despite high baryon density), there is a net pressure differential pushing matter
apart---causing real expansion of the space between galaxy clusters.

This expansion produces real Doppler redshift.\ Light from distant galaxies is redshifted
because those galaxies are genuinely receding, driven apart by the thermal pressure of hot
voids.\ The greater the distance, the more intervening void expansion the light has
traversed, and the greater the cumulative redshift.\ This produces the distance-proportional
redshift described by Hubble's law \cite{hubble1929}.

Critically, because this is real expansion (not an artifact of energy loss or medium
transitions), it naturally produces the time dilation of distant supernova light curves
that is observed and predicted by expansion-based cosmology.\ The light curves are stretched
because the space the light travels through is physically expanding during transit---exactly as in standard cosmological models, but driven by Aether thermal pressure rather
than residual momentum from an initial explosion.


\subsection{Dark Energy}
\label{sec:dark-energy}

The accelerating expansion of the universe finds a natural explanation in this framework.

As matter clusters more tightly under gravity (consuming Aether and locking energy into
mass), it radiates increasing amounts of energy outward into the surrounding voids.\ Stars
emit light, heat, and radiation that propagates into the Aether of deep space.\ Since voids
contain no mass formation to absorb this energy, it accumulates as thermal energy in the
Aether, raising the temperature and thermal pressure of the voids.\ However, the integrated stellar radiation over cosmic history
($\sim$$10^{-14}$~J/m$^3$) falls $\sim$$10^4$ times short of the observed dark
energy density ($\rho_\Lambda c^2 \sim 6 \times 10^{-10}$~J/m$^3$).\ The
two-fluid model (Section~\ref{sec:dark-matter}) provides the missing energy
channel: the normal fraction of the Aether in voids has finite \emph{bulk}
viscosity $\zeta_n$ and dissipates energy from the cosmic expansion.\ The
Hubble flow $\mathbf{v} = H\mathbf{r}$ is pure dilation ($\nabla \cdot
\mathbf{v} = 3H$) with zero shear, so only bulk viscosity contributes.\
The dissipation rate $\dot\varepsilon = \zeta_n\,(3H)^2 = 9\,\zeta_n\,H^2$
accumulates over cosmic time to fill the energy budget, requiring a
normal-component bulk viscosity
%
\begin{equation}
  \zeta_n \;\approx\; \frac{\rho_\Lambda\,c^2}{9\,H^2\,t_0}
    \;\sim\; 3 \times 10^{7}~\text{Pa\,s}
  \label{eq:zeta-normal}
\end{equation}
%
(where $t_0$ is the age of the universe).\ This is an upper bound: $H$ was
larger in the past, so the time-integrated heating $\int 9\,\zeta_n\,H^2\,dt$
exceeds $9\,\zeta_n\,H_0^2\,t_0$ by a factor of $\sim$30--40, giving a more
accurate estimate $\zeta_n \sim 10^6$~Pa\,s.\ This value is an
\emph{observational constraint} on the Aether's microscopic properties: the
observed dark energy density determines what the bulk viscosity must be.\
In superfluids, bulk viscosity arises from relaxation between the superfluid
and normal components (the three-phonon process) and can be substantial
near $T_c$; whether the Aether's excitation spectrum produces
$\zeta_n \sim 10^6$~Pa\,s is a testable prediction that shares the same
unknown---the excitation spectrum---as the dark matter magnitude and the
Hubble tension temperature ratio (Section~\ref{sec:dark-matter},
Section~\ref{sec:hubble-tension}).\ Crucially, photons propagate without dissipation because the superfluid
component has zero shear viscosity, and direct photon--quasiparticle
scattering is suppressed by $P_{\mathrm{th}}/\Ka \sim 10^{-38}$.\
This heating therefore occurs without attenuating electromagnetic
radiation---the two-fluid model provides a natural separation between
bulk-viscous heating (which operates on the normal fraction) and
transparent photon propagation.

This creates a positive feedback loop: matter clusters and radiates energy $\rightarrow$ voids heat
up $\rightarrow$ thermal pressure increases $\rightarrow$ expansion accelerates $\rightarrow$ voids grow larger $\rightarrow$ more
surface area to absorb radiation $\rightarrow$ voids heat up further.

This naturally explains several observed features of dark energy:

\begin{itemize}[leftmargin=2em]
  \item Why dark energy constitutes the vast majority of the universe's energy budget: the
    vast majority of space is void, filled with hot Aether.

  \item Why expansion is accelerating rather than decelerating: the thermal pressure in voids
    is actively increasing as stars continue to radiate energy into them.

  \item Why dark energy appears as a property of space itself rather than a localized force:
    it is a property of the Aether that fills space, driven by the temperature differential
    between voids and matter regions.
\end{itemize}


\subsection{Dark Matter}
\label{sec:dark-matter}

Dark matter---the observation that galaxies and galaxy clusters exhibit stronger
gravitational effects than their visible mass accounts for---is explained in the
Aether framework by the two-fluid first-sound correction: the outward increase
in the first-sound speed $c_1$ from the varying normal fraction provides
additional centripetal acceleration through the acoustic metric, with no
additional mass required.

\paragraph{The two-fluid Aether.}
As a quantum superfluid, the Aether obeys the Landau--Tisza two-fluid model:
below a critical temperature $T_c$, the medium consists of two interpenetrating
components---a superfluid fraction $\rho_s/\rhoa$ (zero viscosity, carries no
entropy) and a normal fraction $\rho_n/\rhoa$ (finite viscosity, carries entropy),
with $\rho_s + \rho_n = \rhoa$ \cite{berezhiani2015}.\ The superfluid fraction
varies continuously with temperature: near matter (cold, $T \ll T_c$),
$\rho_s/\rhoa \approx 1$; in warmer regions (approaching $T_c$), the normal
fraction grows.\ Crucially, photons---spin-wave excitations of the condensate
(Paper~7~\cite{otto-quantum})---propagate without absorption or scattering,
because direct photon--quasiparticle energy exchange is suppressed by
$P_{\mathrm{th}}/\Ka \sim 10^{-38}$.\ The universe therefore remains
transparent regardless of the local normal fraction.\ However, the
normal fraction \emph{does} modify the background propagation speed:
through spin-orbit coupling \cite{volovik2003}, the photon's microscopic
spin-wave character produces macroscopic mass currents indistinguishable
from first-sound perturbations (Paper~7, Section~\ref{sec:qss}), so the
effective photon speed at macroscopic wavelengths is the two-fluid
first-sound speed $c_1$, not the bare condensate speed~$c$.\ The viability of superfluid dark
matter has been independently demonstrated by Berezhiani and Khoury
\cite{berezhiani2015}, who showed that a superfluid medium with appropriate
phonon dynamics can reproduce MOND phenomenology on galactic scales.

\paragraph{Observational reinterpretation.}
Every dark matter measurement ultimately observes a gravitational effect and
infers mass using standard gravity ($v^2 = GM/r$ or $\theta = 4GM/(rc^2)$).\
In the Aether framework, the relation is $v^2 = GM_{\mathrm{visible}}/r +
(r/2)(dc_1^2/dr)$ (Eq.~\ref{eq:vorbit-DM}), so what is conventionally
attributed to ``dark matter mass'' is actually the first-sound gradient
term---no additional mass is required.\ Crucially, both light and matter
follow geodesics of the \emph{same} acoustic metric (light propagates at
$c_1$ through spin-orbit coupling, as described above), so the
$c_1$~gradient that enhances rotation curves automatically produces the
corresponding gravitational lensing.

\paragraph{Quantitative mechanism: two-fluid first-sound correction.}
In the Landau two-fluid model, the first-sound speed---the adiabatic
sound speed $c_1^2 = (\partial P/\partial\rho)_S$---is not simply $c =
\sqrt{\Ka/\rhoa}$ but receives a thermal correction from the normal
component \cite{berezhiani2015}:
%
\begin{equation}
  c_1^2 \;=\; c^2 \;+\; \frac{T\,s^2\,\rho_s}{c_v\,\rho_n}
  \label{eq:first-sound}
\end{equation}
%
where $s$ is the entropy per unit mass, $c_v$ is the specific heat per unit
mass, and $\rho_s/\rho_n$ is the superfluid-to-normal density ratio.\ At
$T = 0$ there are no excitations and trivially $c_1 = c$.\ At any finite
temperature the correction is \emph{positive} ($c_1 > c$); for a pure phonon
gas the correction is $c^2/3$ (temperature-independent), while for a general
excitation spectrum the correction varies with temperature---and it is this
temperature-dependent part that produces the dark matter effect through its
spatial gradient.

\emph{Why photons propagate at $c_1$.}\quad Photons are microscopically
spin-wave excitations of the spinor condensate---transverse oscillations
that do not change the local density (Paper~7~\cite{otto-quantum}).\ However,
through spin-orbit coupling \cite{volovik2003}, these spin oscillations
produce effective mass currents that are indistinguishable from velocity-field
(density/flow) perturbations at macroscopic scales.\ The crossover occurs at
the spin healing length $\xi_s \lesssim 10^{-27}$~m---all observable light is
deeply in the macroscopic regime, where the effective propagation speed is the
first-sound speed $c_1$.\ Transparency is maintained because direct
photon--quasiparticle energy exchange (absorption, scattering) is suppressed
by $P_{\mathrm{th}}/\Ka \sim 10^{-38}$---a measure of scattering cross-section,
not of the background speed modification.\ The $c_1$ correction is purely
thermodynamic and therefore wavelength-independent: gravitational lensing from
the dark matter component is achromatic, consistent with observations.

Since the normal fraction increases outward from the galaxy (warmer Aether has
more thermal excitations), $c_1$ \emph{increases outward}.\ Extending the
acoustic metric of Paper~3 (Section~\ref{sec:lensing}) to position-dependent
wave speed $c_1(r)$, the circular-orbit condition $v_{\mathrm{orbit}}^2 =
(r/2)\,dg_{tt}/dr$ gives:
%
\begin{equation}
  v_{\mathrm{orbit}}^2 \;=\; \frac{GM}{r}
    \;+\; \frac{r}{2}\,\frac{dc_1^2}{dr}
  \label{eq:vorbit-DM}
\end{equation}
%
the positive gradient $dc_1^2/dr > 0$ provides additional centripetal
acceleration---precisely the ``dark matter'' effect.\ The sign is
\emph{automatically correct}: it follows from the Landau two-fluid
thermodynamics without any assumptions about which phase is stiffer.\ The
effect is naturally cored: in the pure superfluid center ($\rho_n \approx 0$),
the correction vanishes and only visible-mass gravity operates.

\paragraph{Magnitude of the first-sound correction.}
The dark matter effect depends on the \emph{spatial variation}
$\Delta c_1^2$ across the halo, which requires the correction to vary
with temperature.\ The magnitude of this variation depends critically on
the Aether's entropy per unit mass $s$, which in turn depends on the
full excitation spectrum (phonons, and any roton-like or other modes
beyond the long-wavelength limit).\ Two limiting cases bracket the answer:
%
\begin{itemize}[leftmargin=2em]
  \item \textbf{Phonon gas only} (long-wavelength collective modes):
    For a pure phonon gas ($\omega = ck$), the correction $Ts^2\rho_s/(c_v\rho_n) = c^2/3$
    is \emph{temperature-independent} (the phonon scaling $s \propto T^3$,
    $\rho_n \propto T^4$ produces exact cancellation).\ A constant correction
    has zero spatial gradient and produces no dark matter effect; the residual
    variation from dispersion corrections is
    $\Delta c_1^2/c^2 \sim 10^{-14}$---negligibly small.
  \item \textbf{Ideal BEC} (all single-particle excitations):
    $\Delta c_1^2/c^2 \sim 10^{-4}$---$\sim$$50\times$ larger than
    needed.
\end{itemize}
%
The observed dark matter requires $\Delta c_1^2/c^2 \sim 10^{-6}$ across
the halo, corresponding to an entropy intermediate between these limits.\
For comparison, helium-4 at $T/T_c \approx 0.7$ has entropy dominated by
roton excitations, placing it near the ideal-BEC limit; whether the Aether
has analogous intermediate-wavelength excitations is the key open question.

\paragraph{Quantitative consistency across observations.}
A flat rotation curve ($v_{\mathrm{flat}} = \mathrm{const}$) requires
$dc_1^2/dr = 2v_{\mathrm{flat}}^2/r$, giving a logarithmic profile
$\Delta c_1^2(r) = 2v_{\mathrm{flat}}^2\ln(r/r_0)$.\ The total variation
across a Milky-Way-sized halo ($v_{\mathrm{flat}} = 220$~km/s, 1--200~kpc)
is $\Delta c_1^2/c^2 \approx 6 \times 10^{-6}$---six parts per million,
undetectable by any terrestrial experiment.

\emph{Gravitational lensing.}\ Because both light and matter follow
geodesics of the same acoustic metric, the effective mass profile from
the $c_1$~gradient is exactly the singular isothermal sphere (SIS):
$M_{\mathrm{eff}}(r) = v_{\mathrm{flat}}^2 r/G$.\ The corresponding
deflection angle is $\theta = 2\pi v_{\mathrm{flat}}^2/c^2 \approx
0.70''$, identical to the SIS result $\theta_{\mathrm{SIS}} = 4\pi
\sigma^2/c^2$ (with $\sigma = v_{\mathrm{flat}}/\sqrt{2}$).\ No separate
lensing calculation is needed---rotation curves and lensing are
automatically consistent.

\emph{Galaxy clusters.}\ For a massive cluster ($\sigma \sim
1000$~km/s), the required $\Delta c_1^2/c^2 \approx 10^{-4}$---roughly
$17\times$ larger than for a galaxy.\ This is physically expected: clusters
are hotter (more normal fraction, higher entropy), so the first-sound
correction is naturally larger.

\emph{Bullet Cluster.}\ The $c_1$ correction is maintained by each
galaxy's cooling of the surrounding Aether through the binding process.\
During the collision, the shocked gas heats the local Aether, reducing the
normal-fraction contrast and weakening the $c_1$ correction in the collision
region.\ The surviving corrections follow the galaxies, reproducing the
observed lensing--gas separation.

\emph{No dark matter particles.}\ The mechanism is a thermodynamic
property of the two-fluid medium, not a new particle species, explaining
the null results of direct-detection experiments.

\emph{Tully--Fisher relation.}\ The observed baryonic relation
$M \propto v^4$ constrains the model.\ For steady-state thermal
conduction, $T_{\mathrm{void}} - T(R_{25}) = L_{\mathrm{cool}}/(4\pi
\kappa R_{25})$, and the orbital velocity at $R_{25}$ is
$v^2 \propto (L_{\mathrm{cool}}/R_{25})^{1+\beta}$, where
$\beta$ parametrizes the critical behavior of the first-sound
correction near $T_c$: $f'(T) \propto (T_{\mathrm{void}} - T)^\beta$.\
The Tully--Fisher slope depends on \emph{both} $\beta$ and how the
galaxy's cooling rate scales with mass ($L_{\mathrm{cool}} \propto
M^\alpha$).\ For $\alpha = 1$ (linear cooling) and $\beta = 0$
(smooth, ideal BEC), the model gives $M \propto v^3$ (wrong slope).\
However, $\beta = -1/4$ (a weak cusp in $f'(T)$, as seen in
helium-4's lambda anomaly) yields $M \propto v^4$ exactly.\
Alternatively, sublinear cooling ($\alpha \approx 0.7$, as expected
from the observed star-formation main sequence) with smooth $\beta
\approx +0.35$ also yields $M \propto v^4$.\ The correct slope thus
depends on the same excitation spectrum and critical behavior that
determine the dark matter magnitude---an open but specific calculation.

\paragraph{Quantitative status and open calculations.}
The first-sound correction mechanism has several structural advantages:
(a)~the sign is correct by construction (Landau thermodynamics);
(b)~cored profiles emerge naturally (no correction in the pure superfluid
center); (c)~it is independent of the quantum of circulation $\kappa_0$;
(d)~rotation curves and lensing are automatically consistent (same
acoustic metric); (e)~the scaling from galaxies to clusters is
qualitatively correct.\ The critical open calculations are:
(i)~the Aether's excitation spectrum beyond the phonon regime, which
determines the entropy and hence the magnitude ($\Delta c_1^2/c^2$ must
be $\sim$$10^{-6}$ for galaxies; the phonon-only limit gives
$10^{-14}$ and the ideal BEC gives $10^{-4}$);
(ii)~the three-dimensional temperature profile $T(r,z)$ from first
principles; (iii)~the Tully--Fisher exponent (currently $\sim$3 vs.\
observed $\sim$4); (iv)~the CMB power spectrum from the cosmological
superfluid phase transition; (v)~whether the $c_1$ correction propagates
differently from collisionless dark matter during acoustic oscillations.\
The viability of superfluid dark matter has been demonstrated in
principle by Berezhiani and Khoury \cite{berezhiani2015} through
phonon-mediated forces; the Aether's two-fluid first-sound mechanism
operates through a different channel but within the same superfluid
dark matter paradigm.


\subsection{Cosmic Microwave Background Radiation}

In this theory, photons are spin-wave excitations of the Aether condensate
(Paper~7~\cite{otto-quantum}).\ The cosmic microwave background is therefore
not radiation traveling \emph{through} the medium but the thermal excitation
spectrum \emph{of} the medium---the Aether at thermal equilibrium.\ The CMB
temperature ($T_{\mathrm{CMB}} = 2.725$~K) is the Aether's temperature, and
the near-perfect blackbody spectrum confirmed by COBE/FIRAS \cite{fixsen1996}
follows naturally: it is the equilibrium spectrum of a quantum medium.

The CMB anisotropies ($\delta T/T \sim 10^{-5}$) reflect spatial temperature
variations in the Aether.\ The specific angular power spectrum---with its
characteristic acoustic peaks encoding baryon--photon physics---remains an open
question within the Aether framework.\ Two candidate mechanisms are identified
in the Open Questions section: the Kibble--Zurek defect spectrum from the
cosmological superfluid phase transition, and second-sound oscillations in the
two-fluid Aether.\ Deriving the CMB power spectrum from first principles is one
of the critical open calculations for this theory.

This theory does not preclude the Big Bang; it offers complementary mechanisms
for phenomena traditionally attributed solely to it.


\subsection{Observational Evidence for Dynamic Dark Energy}

The thermal pressure mechanism makes a distinctive prediction: dark energy should
\emph{evolve over time}.\ As voids expand and accumulate radiative energy, their
thermal pressure changes---it is not a fixed cosmological constant but a dynamic
quantity driven by the ongoing thermodynamic interplay between matter-containing
regions and voids.\ This contrasts sharply with the standard $\Lambda$CDM model,
in which the cosmological constant $\Lambda$ is strictly time-independent.

Recent observations strongly favor the dynamic prediction.\ The Dark Energy
Spectroscopic Instrument (DESI) collaboration, using baryon acoustic oscillation
measurements from over 14 million galaxies and quasars \cite{desi2025}, has found
evidence at 2.8--4.2$\sigma$ significance that the dark energy equation of state
parameter $w$ varies with redshift.\ The data prefer a model in which the dark energy
density is \emph{slowly decreasing} with time, inconsistent with a cosmological
constant at the $\sim$3$\sigma$ level when combined with cosmic microwave background
and supernovae data.

This is precisely what the Aether thermal pressure model predicts.\ Two competing
effects govern the evolution:
%
\begin{itemize}[leftmargin=2em]
  \item \textbf{Heating}: Stars continuously radiate energy into voids, raising
    the thermal pressure.\ This drives the positive feedback loop described above.

  \item \textbf{Dilution}: As voids expand, the Aether within them becomes less
    dense.\ Since pressure depends on both temperature and density, the net effect
    on expansion rate depends on the balance between heating and dilution.
\end{itemize}
%
The observed decrease in dark energy density is consistent with dilution beginning
to dominate over heating at late cosmological times---a natural outcome when the
rate of void expansion outpaces the rate of energy injection from matter regions.

An independent line of evidence comes from the Finsler gravity framework
\cite{hohmann2026}, which extends general relativity to incorporate the thermal
pressure of kinetic gases directly into the field equations.\ When this is done,
the Friedmann equations naturally produce cosmic acceleration without requiring
a separate dark energy term---the acceleration arises from the thermal pressure
of the medium itself.\ While the mathematical formalism differs from the Aether
model, the physical mechanism is strikingly parallel: in both cases, cosmic
acceleration is driven by the thermal pressure of a space-filling medium rather
than by a cosmological constant.

The Dark Energy Survey's final analysis \cite{des2026}, combining weak lensing and
galaxy distribution measurements over 6 billion years of cosmic history, provides
the tightest constraints yet on the dark energy equation of state---and these
constraints are consistent with a slowly varying, dynamic dark energy rather than
a strict cosmological constant.


\subsection{Void Structure and Differential Expansion}

The Aether framework predicts that cosmic voids should exhibit specific, universal
properties dictated by the interplay of density, pressure, and temperature.\ In
particular, since $\Ka$ is invariant and $\rhoa$ varies smoothly, the void density
profile should be governed by the same physics everywhere---producing similar
profiles across different void sizes and cosmic epochs.

Galaxy surveys confirm this prediction.\ When void density profiles are scaled by
the void effective radius, the resulting profiles are approximately universal:
the shape is largely independent of void size and redshift.\ This universality
is difficult to explain if voids are shaped purely by initial conditions and
gravitational dynamics (which would produce diverse profiles depending on local
environment and formation history), but follows naturally from a medium whose
bulk modulus $\Ka$ is the same everywhere.

Furthermore, galaxies within voids show systematically different properties from
those in denser environments: lower stellar masses, bluer colors, later
morphological types, higher gas fractions, and lower metallicities.\ These
``delayed evolution'' signatures are consistent with the higher Aether
temperature in voids: the stronger first-sound correction (larger normal
fraction) shifts the effective gravitational dynamics, and the higher
thermal energy potentially raises thresholds for permanent binding (the
DCE threshold depends on local conditions).

The timescape cosmology program \cite{wiltshire2025} has demonstrated that
differential expansion rates between voids and walls---exactly the behavior
predicted by the Aether three-property framework---can account for the observed
distance-redshift relation of Type Ia supernovae without requiring dark energy
at all.\ In that model, as in the Aether model, voids expand faster than
matter-containing regions because they contain less gravitational deceleration.
The Aether framework adds a physical mechanism: the thermal pressure differential
between hot voids and cold matter regions actively drives the differential
expansion.


\subsection{The Hubble Tension}
\label{sec:hubble-tension}

One of the most pressing open problems in modern cosmology is the \emph{Hubble
tension}: a persistent, statistically significant discrepancy between
early-universe and late-universe measurements of the expansion rate $H_0$.
Observations of the cosmic microwave background by the Planck satellite, analyzed
within $\Lambda$CDM, yield $H_0 = 67.4 \pm 0.5$~km\,s$^{-1}$\,Mpc$^{-1}$
\cite{planck2020}.\ Local distance-ladder measurements using Cepheid-calibrated
Type~Ia supernovae give $H_0 = 73.04 \pm 1.04$~km\,s$^{-1}$\,Mpc$^{-1}$
\cite{riess2022}---a discrepancy of $\sim$8\%, exceeding 5$\sigma$ significance.
Independent measurements from gravitational lensing time delays also favor the
higher value \cite{wong2020}, and the James Webb Space Telescope has ruled out
Cepheid crowding artifacts at $>$8$\sigma$ confidence.\ The tension is
increasingly regarded as evidence that the standard cosmological model is
incomplete.

The Aether framework addresses this tension through spatially varying dark
energy: because the thermal pressure that drives cosmic expansion depends on
local temperature, different environments expand at different rates.

\paragraph{The local void.}
Observational surveys have established that our galaxy resides within the
Keenan--Barger--Cowie (KBC) void \cite{haslbauer2020}, a region approximately
300~Mpc in radius.\ The observed luminosity density contrast is
$\delta \equiv 1 - \rho/\bar\rho \approx 0.46 \pm 0.06$, indicating that the
local universe is roughly 46\% less dense in \emph{matter} than the cosmic
average.\ Approximately half of this apparent contrast is an artifact of
void-induced outflows that distort the redshift--distance relation; the physical
matter underdensity is estimated at $\sim$20--30\%.

In $\Lambda$CDM, the KBC void poses a severe challenge: a void of this size
and depth is expected fewer than 1 in 10$^{9}$ times in standard simulations
\cite{haslbauer2020}.\ The Aether framework resolves this problem.\ The thermal
pressure differential between hot void interiors and cold matter-containing
walls (Section~1.2) actively creates and sustains large voids---the KBC void is
not a statistical fluke but the natural outcome of the Aether's three-property
dynamics.

\paragraph{Why the standard void outflow alone does not resolve the tension.}
In an underdense region, matter experiences a net outward peculiar velocity
driven by the gravitational pull of the denser surrounding walls.\ In linear
perturbation theory, the peculiar velocity within a spherical underdensity
of contrast $\delta$ is:
%
\begin{equation}
  v_{\mathrm{pec}}(r) \;\approx\; -\frac{1}{3}\,H_0\,f(\Omega_m)\,\delta\,r
  \label{eq:vpec}
\end{equation}
%
where $f \approx \Omega_m^{0.55} \approx 0.53$.\ For a physical underdensity of
$\sim$25\%, this outflow contributes $\Delta H_0 \approx +3$~km\,s$^{-1}$\,Mpc$^{-1}$.\ However, the void's matter deficit also \emph{reduces}
the local expansion rate through the Friedmann equation: with less matter
contributing to the total energy density, $H_{\mathrm{local}}^2 \propto
\Omega_m(1 - \delta) + \Omega_\Lambda$ is lower than the cosmic average,
contributing $\sim$$-$2.6~km\,s$^{-1}$\,Mpc$^{-1}$.\ These two effects nearly
cancel, leaving a net boost of only $\sim$0.4~km\,s$^{-1}$\,Mpc$^{-1}$---far
short of the observed $+$5.6~km\,s$^{-1}$\,Mpc$^{-1}$ discrepancy.\ This
cancellation is well established: analyses within standard $\Lambda$CDM have shown
that local voids cannot resolve the Hubble tension when dark energy is spatially
uniform.

\paragraph{Resolution through temperature-dependent dark energy.}
The Aether model breaks this cancellation.\ In $\Lambda$CDM, the cosmological
constant $\Lambda$ is spatially uniform---voids and walls experience the same
dark energy density.\ In the Aether model, dark energy is the thermal pressure
of the Aether, which depends on local \emph{temperature} (Section~1.2).\
Crucially, the Aether density $\rhoa$ is essentially uniform across cosmological
structures: because $\Ka \sim 10^{28}$~J/m$^3$ is enormous, the Aether resists
compression and rarefaction far more effectively than any cosmological force can
drive it.\ The gravitational density perturbation from the KBC void is only
$\delta\rhoa/\rhoa \sim 2G\Delta M/(Rc^2) \sim 10^{-4}$, so the speed of
light $c = \sqrt{\Ka/\rhoa}$ is effectively constant across cosmic structures
(varying by less than 0.01\%).\ Similarly, although $G \propto \epsilon\,\rhoa\,\psi^2$ is in
principle spatially variable (Section~1.3), both channels of variation are
negligible here: the density channel contributes $\delta G/G \sim
\delta\rhoa/\rhoa \sim 10^{-4}$, and the temperature channel (through the
asymmetry fraction $\epsilon$) is suppressed by the ratio
$P_{\mathrm{th}}/\Ka \sim 10^{-37}$---the thermal energy is far too small
relative to the medium's mechanical energy scale to perturb the bind--unbind
cycle.\ The Friedmann equation below therefore uses constant~$G$.\
(Intriguingly, at the CMB temperature this ratio drops to $\sim\!10^{-43}$,
numerically equal to the asymmetry fraction~$\epsilon$ itself;
Paper~3~\cite{otto-gravity} explores the speculative possibility that the two
are physically identified through the normal-fluid fraction of the two-fluid
model, yielding $T \approx T_{\mathrm{CMB}}$ for
$\rhoa \approx 1.3 \times 10^{12}$~kg/m$^3$.)

Temperature, by contrast, varies dramatically.\ Voids accumulate radiative
energy with no mass formation to absorb it, reaching temperatures far above
matter-containing regions.\ Since the thermal pressure depends on temperature,
the effective dark energy is higher in voids and lower near matter.

This modifies the local Friedmann equation.\ Define $\eta$ as the fractional
thermal pressure excess of the void above the cosmic average, so that the
effective dark energy density in the void is $\Omega_\Lambda(1 + \eta)$.\ The
locally measured $H_0$ then has two contributions:
%
\begin{equation}
  H_{0,\mathrm{local}} \;=\; \underbrace{H_{0,\mathrm{avg}}\,
    \sqrt{\,\Omega_m(1 - \delta) + \Omega_\Lambda(1 + \eta)\,}}_{\text{local
    Friedmann rate}} \;+\;
    \underbrace{\tfrac{1}{3}\,H_{0,\mathrm{avg}}\,f\,\delta}_{\text{peculiar
    velocity}}
  \label{eq:H0-local}
\end{equation}
%
The first term is the local expansion rate (modified by both the matter
deficit and the thermal pressure excess); the second is the gravitational
outflow from the void density gradient.

For $\Omega_m = 0.3$, $\Omega_\Lambda = 0.7$, and $\delta = 0.25$, matching the
observed $H_{0,\mathrm{local}} = 73.0$~km\,s$^{-1}$\,Mpc$^{-1}$ requires:
%
\begin{equation}
  \eta \;\approx\; 0.22
  \label{eq:eta-required}
\end{equation}
%
That is, the KBC void's thermal pressure must be approximately 22\% above the
cosmic average.\ This is the analysis's single fitted parameter: $\Omega_m$,
$\Omega_\Lambda$, $\delta$, $f$, and $H_{0,\mathrm{avg}}$ are all externally
measured, and the peculiar velocity and Friedmann deficit are derived from
standard physics with no free parameters; only $\eta$ is determined by
matching the observed local expansion rate.\ The result is robust: varying
$\delta$ across the plausible range 0.15--0.30 changes $\eta$ by less than
$\pm$0.01.

If voids occupy a fraction $f_v \approx 0.7$ of the cosmic volume and thermal
energy is conserved, the void and wall temperatures satisfy:
%
\begin{equation}
  \frac{T_{\mathrm{void}}}{T_{\mathrm{wall}}} \;\approx\; 2.5
  \label{eq:T-ratio}
\end{equation}
%
This ratio assumes the thermal pressure is proportional to temperature
($P_{\mathrm{th}} \propto T$, as for an ideal gas); a radiation-dominated
pressure ($P \propto T^4$) would give a smaller ratio of $\sim$1.3.\ A
factor-of-2.5 temperature contrast is physically plausible: voids accumulate
radiative energy over cosmological timescales with no mass formation to absorb
it, while matter-containing walls continuously lock thermal energy into mass
through the binding mechanism (Section~1.2).\ Whether the Aether's thermodynamic
history quantitatively produces this ratio is an open question (see below).\
Notably, the pressure law ($P \propto T^n$) depends on the Aether's excitation
spectrum---the same unknown that determines the dark matter magnitude
(Section~\ref{sec:dark-matter}), the bulk viscosity $\zeta_n$
(Section~\ref{sec:dark-energy}), and the dark energy equation of state
$w(z)$.\ This is a powerful consistency test: a single excitation spectrum
must simultaneously produce the correct dark matter halo profiles
($\Delta c_1^2/c^2 \sim 10^{-6}$), the correct bulk viscosity for the dark
energy density ($\zeta_n \sim 10^6$~Pa\,s), the correct void temperature
ratio for the Hubble tension ($T_{\mathrm{void}}/T_{\mathrm{wall}} \approx
2.5$), and the correct $w(z)$ evolution for DESI.\ If any one of these
constraints is incompatible with the others, the theory fails.

The decomposition of effects illustrates why the Aether-specific contribution
is essential:
%
\begin{center}
\begin{tabular}{@{} l r l @{}}
\toprule
Effect & $\Delta H_0$ & Source \\
\midrule
Matter deficit (Friedmann)  & $-2.6$ & Less matter $\rightarrow$ lower $H^2$ \\
Peculiar velocity outflow   & $+3.0$ & Gravitational outflow from void \\
\emph{Net without Aether}   & $+0.4$ & \emph{Nearly cancels} \\[4pt]
Thermal pressure excess ($\eta \approx 0.22$) & $+5.2$ & Hotter void $\rightarrow$ higher $H^2$ \\
\midrule
\textbf{Total}              & $+5.6$ & Matches observed tension \\
\bottomrule
\end{tabular}
\end{center}
%
The thermal pressure enhancement does essentially all of the work.\ The standard
void outflow ($+$3.0~km\,s$^{-1}$\,Mpc$^{-1}$) is almost exactly canceled by the
Friedmann deficit ($-$2.6~km\,s$^{-1}$\,Mpc$^{-1}$), leaving the spatially
varying dark energy as the dominant mechanism.\ The timescape cosmology program
\cite{wiltshire2025}, which demonstrates that differential expansion between
voids and walls can account for Type~Ia supernova data, provides independent
support for this class of environment-dependent expansion.

\paragraph{Evolving dark energy: a confirmed prediction that does not directly
resolve the tension.}
The DESI collaboration has found evidence at 2.8--4.2$\sigma$ significance
\cite{desi2025} that the dark energy equation of state evolves with redshift,
with best-fit parameters indicating a phantom-to-quintessence crossing around
$z \approx 0.5$.\ The Aether thermal pressure model predicts precisely this
behavior: at early times, voids were smaller and denser with high thermal
pressure relative to their volume (effective $w < -1$); at late times, dilution
of expanding void Aether dominates over heating from stellar radiation, producing
a decreasing dark energy density ($w > -1$).

This is an important confirmation of the Aether dark energy model.\ However, the
specific phantom crossing favored by DESI does \emph{not} help resolve the
Hubble tension.\ When DESI and Planck data are jointly fit with $w_0 w_a$CDM,
the inferred $H_0$ shifts by less than 0.3~km\,s$^{-1}$\,Mpc$^{-1}$ relative
to $\Lambda$CDM---leaving the tension essentially unchanged.\ The physical
reason is that the phantom-to-quintessence transition implies decreasing dark
energy density at low redshift, which \emph{reduces} late-time acceleration.

In the Aether framework, this is consistent: the Hubble tension arises from
\emph{local} temperature variations (the void is hotter), not from the
\emph{global} expansion history.\ The evolving dark energy and the local void
are independent features of the same three-property Aether model, and the DESI
confirmation validates the dark energy prediction on its own merits.

\paragraph{Testable predictions.}
The Aether framework for the Hubble tension makes several predictions that
distinguish it from other proposed solutions:
%
\begin{enumerate}[leftmargin=2em]
  \item \textbf{Directional anisotropy}: Since the KBC void is not perfectly
    spherical, the locally measured $H_0$ should vary across the sky.\ The Aether
    model predicts the anisotropy pattern traces the \emph{temperature} profile
    of the void (which determines the thermal pressure), not merely the matter
    density profile.

  \item \textbf{Void temperature--expansion correlation}: If thermal pressure
    drives the differential expansion, there should be a measurable correlation
    between the thermal energy content of individual voids (inferred from their
    ISW imprints or SZ signals) and the local expansion rate measured within or
    behind those voids.

  \item \textbf{Convergence scale}: The tension should diminish as measurements
    average over larger volumes that encompass both voids and walls.\ The Aether
    model predicts convergence at the scale of statistical homogeneity
    ($\sim$100~Mpc), with a specific functional form governed by the void
    temperature distribution.

  \item \textbf{Void abundance}: The Aether model predicts that large voids are
    more common than $\Lambda$CDM simulations allow, because thermal pressure
    actively sustains them.\ Upcoming surveys (DESI, Euclid) will map the void
    size distribution with sufficient precision to test whether $\Lambda$CDM's
    predicted void abundance matches observations.
\end{enumerate}


\subsection{The CMB Cold Spot and Supervoid Anomalies}

The CMB Cold Spot---a $\sim$5$^\circ$ radius feature approximately 150~$\mu$K
colder than the mean CMB temperature in the direction of the constellation
Eridanus---has remained unexplained since its detection \cite{vielva2004}.\ A
supervoid approximately 1.8 billion light-years across has been confirmed in
the same direction \cite{szapudi2015}, but under $\Lambda$CDM, the integrated
Sachs-Wolfe (ISW) effect from this void accounts for only 10--20\% of the
observed temperature depression.

More broadly, surveys of supervoids elsewhere in the sky have found ISW
temperature imprints approximately $5.2 \pm 1.6$ times stronger than $\Lambda$CDM
predicts \cite{kovacs2022}---a significant anomaly that suggests the standard
model underestimates the influence of voids on CMB photons.

The Aether framework offers a potential resolution.\ In standard cosmology, the
ISW effect depends only on the time derivative of the gravitational potential as
photons traverse a void.\ In the Aether model, photons traversing a void also
encounter different medium properties:
%
\begin{itemize}[leftmargin=2em]
  \item \textbf{Modified flow environment}: Photons traversing a void encounter an
    outward-expanding Aether flow (driven by the thermal pressure excess
    $\eta \approx 0.22$; Section~1.7), which modifies their energy through the
    acoustic metric differently than a uniform $\Lambda$ would.

  \item \textbf{Higher temperature / lower $\epsilon$}: The hot void Aether has a
    lower gravitational asymmetry fraction $\epsilon$, weakening the gravitational
    potential more rapidly as the void expands---amplifying the ISW time derivative
    $\partial\Phi/\partial t$.

  \item \textbf{Expansion asymmetry}: The void expands faster than the
    surrounding regions (enhanced thermal pressure), so photons lose more energy
    climbing out of the expanding void than they gained falling in---but the
    magnitude depends on the Aether temperature profile, not just the
    gravitational potential.
\end{itemize}
%
The net effect could amplify the ISW signal beyond the $\Lambda$CDM prediction,
potentially explaining the factor-of-five discrepancy observed in supervoid
temperature imprints.\ The thermal pressure excess $\eta \approx 0.22$ from the
Hubble tension analysis (Eq.~\ref{eq:eta-required}) provides a concrete starting
point: the modified Friedmann driving term amplifies $\partial\Phi/\partial t$
by a factor of $\sim$1.6, with the temperature-dependent $\epsilon$ and
expansion asymmetry contributing additional amplification.\ A full quantitative
calculation remains to be completed (see Open Questions).


% ============================================================
\section{Black Holes}
\label{sec:black-holes}
% ============================================================

The event horizon, derived in the gravitational lensing section as a sonic horizon where
the Aether inflow velocity reaches $c$, provides the foundation for a complete Aether-based
treatment of black hole physics.\ Several observed and predicted properties of black holes
emerge naturally from the inflow model.


\subsection{Hawking Radiation}

In 1974, Hawking \cite{hawking1974,hawking1975} predicted that black holes are not perfectly black---they emit thermal
radiation at a temperature inversely proportional to their mass:
%
\begin{equation}
  T_{\mathrm{Hawking}} = \frac{\hbar \, c^3}{8\pi \, G \, M \, k_B}
\end{equation}
%
This prediction has not been directly observed (the radiation is far too weak for stellar-mass black holes), but it is widely accepted on theoretical grounds and has been confirmed
in acoustic analogue systems \cite{steinhauer2019}.

In the Aether model, Hawking radiation has a concrete physical mechanism rooted in the
transient binding cycle.

Far from any mass, the binding cycle is nearly symmetric: Aether binds into a transient
pair, the pair annihilates, and the restoring pressure pulse propagates outward.\ The net
energy loss is only the tiny asymmetry $\epsilon$---the gravitational wave.

Near the event horizon, this symmetry breaks down dramatically.\ The Aether is flowing
inward at nearly $c$, so the restoring pressure pulse from pair annihilation---which
propagates outward at the local speed of light---is severely impeded by the inflow.\ The
outward-propagating pulse is redshifted, weakened, and partially swept back inward.\ The
fraction of each cycle's energy that escapes outward is no longer the tiny $\epsilon$ of
normal-space gravity, but a significantly larger fraction determined by the inflow
velocity gradient at the horizon.

Just outside the horizon, some of each cycle's restoring energy still escapes.\ Just
inside, none does---the inflow exceeds $c$ and all restoring pulses are swept inward.\ The
sharp transition at the horizon means that transient binding events straddling the horizon
produce a net outward energy flux: the portion of the restoring pulse that forms outside
the horizon escapes, while the portion that forms inside is lost.\ This asymmetry is the
physical mechanism of Hawking radiation.

The temperature of this radiation depends on the gradient of the inflow velocity at the
horizon---steeper gradients produce sharper asymmetry and hotter radiation.\ The inflow
velocity gradient at the Schwarzschild radius is:
%
\begin{equation}
  \left|\frac{d\vin}{dr}\right|_{r_s} = \frac{c^3}{4GM}
\end{equation}
%
This is proportional to $1/M$: smaller black holes have steeper gradients and therefore
higher Hawking temperatures.\ This is consistent with Hawking's prediction that
$T \propto 1/M$.

The Aether mechanism also connects Hawking radiation to the Unruh effect (identified as
an open question elsewhere in this paper).\ An observer accelerating through Aether at
acceleration $a$ experiences the same binding asymmetry as an observer hovering near a
horizon with surface gravity $a$---the acceleration creates an effective inflow gradient
that disrupts the symmetry of the local binding cycle.\ In the Aether model, the Hawking
effect and the Unruh effect are the same phenomenon: asymmetric transient binding caused
by velocity gradients in the medium, whether those gradients arise from gravitational
inflow or from the observer's acceleration.

This identification has a significant theoretical consequence.\ The acoustic analogue of
Hawking radiation has been experimentally confirmed in Bose--Einstein condensate systems
\cite{steinhauer2019}, where a flowing superfluid creates a sonic horizon and emits thermal
phonons at the predicted temperature.\ In the Aether model, this is not merely an analogy---it is the same physical process occurring in a different medium.\ The BEC experiment
confirms that sonic horizons in flowing media produce thermal radiation, which is exactly
what the Aether model predicts for gravitational event horizons.


\subsection{The Photon Sphere and Black Hole Shadow}

In general relativity, photons can orbit a non-rotating black hole at a radius of
$r_{\mathrm{photon}} = 3GM/c^2$, equal to 1.5 times the Schwarzschild radius.\ This photon sphere
defines the apparent size of the black hole shadow---the dark silhouette observed by
the Event Horizon Telescope in its images of M87* \cite{eht2019} and Sagittarius~A* \cite{eht2022}.

In the Aether model, the photon sphere arises from the same two mechanisms that produce
gravitational lensing: the acoustic metric of the Aether inflow (transverse
advection plus time-delay asymmetry).\ At the photon sphere radius, the combined
inward bending from these two
effects is strong enough to bend light into a circular orbit.\ The weak-field bending
formula $4GM/(r_0 c^2)$ applies to rays passing far from the mass; near the photon
sphere, the strong-field regime requires the full density and inflow profiles to be solved
self-consistently.

The Event Horizon Telescope measurements of the M87* shadow are consistent with the
general-relativistic prediction to approximately 10\% precision.\ Since the Aether model
reproduces the weak-field bending angle exactly ($4GM/(r_0 c^2)$), the strong-field
extension is expected to yield the correct photon sphere radius and shadow size, but this
requires quantitative verification through numerical ray-tracing in the Aether density
and inflow profiles.


\subsection{Frame Dragging}

A rotating mass drags the surrounding spacetime into co-rotation---an effect predicted
by general relativity (the Lense--Thirring effect \cite{lense1918}) and confirmed by the Gravity Probe~B
experiment \cite{everitt2011} (2011), which measured the frame-dragging precession of gyroscopes in Earth
orbit to be $37.2 \pm 7.2$ milliarcseconds per year, consistent with the predicted
$39.2$ milliarcseconds per year.

The quantitative derivation of frame dragging from the $v^2/c^2$
universality is given in Section~\ref{sec:gravitomagnetism}.

In the Aether model, frame dragging is a natural consequence of the vortex mechanism.\ Each
electron produces a spin vortex that circulates Aether around its spin axis.\ In a rotating
body, the aggregate vortex circulation of all its electrons produces a net azimuthal Aether
flow around the body.\ This is directly analogous to a rotating body in a viscous fluid
dragging the fluid into co-rotation---except that the coupling mechanism is the
electromagnetic vortex interaction rather than viscous friction.

For a rotating black hole, the frame-dragging effect creates an ergosphere---a region
outside the event horizon where the Aether is dragged so rapidly that no object can remain
stationary relative to the distant Aether.\ Within the ergosphere, all matter and light
must co-rotate with the inflowing, rotating Aether.\ The Penrose process \cite{penrose1969} (extracting energy
from a rotating black hole) has a natural interpretation: an object entering the ergosphere
can exchange angular momentum with the rotating Aether flow, gaining energy at the expense
of the black hole's rotational kinetic energy.


\subsection{The No-Hair Theorem}

In general relativity, black holes are characterized by only three externally measurable
quantities: mass, angular momentum (spin), and electric charge.\ All other information
about the matter that formed the black hole is lost behind the event horizon.

The Aether model provides a physical explanation for this result.\ Inside the event
horizon, the Aether inflow exceeds $c$, sweeping all internal structure inward.\ The only
properties that can influence the external Aether flow pattern are:

\begin{itemize}[leftmargin=2em]
  \item \textbf{Mass:} determines the inflow velocity profile $\vin(r) = \sqrt{2GM/r}$ and therefore the
    inflow velocity profile, time dilation, and gravitational wave output.

  \item \textbf{Angular momentum:} determines the azimuthal Aether flow pattern (frame dragging).
    The net spin of the infalling matter is preserved as rotational motion of the Aether
    inflow.

  \item \textbf{Charge:} determines the net radial Aether flow direction and magnitude (electric
    field).\ A net positive or negative charge produces a persistent radial outflow or
    inflow superimposed on the gravitational inflow.
\end{itemize}

All other properties of the infalling matter---its composition, temperature, shape,
internal structure---are swept behind the sonic horizon and cannot influence the external
Aether flow.\ The no-hair theorem is a consequence of the sonic horizon's one-way nature:
only quantities that produce long-range Aether flow patterns (gravity, rotation,
electromagnetic field) survive.


\subsection{Interior Structure and the Singularity Question}

General relativity predicts that all matter inside a black hole collapses to a singularity---a point of infinite density.\ This prediction is widely regarded as a sign that GR
breaks down at extreme scales, requiring a theory of quantum gravity for resolution.

The Aether model offers a different perspective.\ Inside the event horizon, the Aether
inflow exceeds $c$, and the restoring pressure pulses from transient binding can no longer
propagate outward.\ This has several consequences:

\begin{itemize}[leftmargin=2em]
  \item The transient binding cycle continues (pair formation and annihilation are local
    processes), but the gravitational wave energy that would normally escape as the tiny
    asymmetry $\epsilon$ is instead carried further inward by the flow.\ This represents a net
    inward energy transport that does not occur outside the horizon.

  \item As Aether accumulates toward the center, the density increases to extreme values.\ At
    some density, the perfect linearity of the Aether's pressure--density relationship
    ($\Ka = \text{constant}$) may break down.\ If $\Ka$ increases with extreme compression (the
    medium stiffens), the inflow would encounter increasing resistance, potentially
    reaching an equilibrium at finite density rather than collapsing to infinite density.

  \item At sufficiently extreme density, the binding threshold may be trivially exceeded
    everywhere, converting all Aether to bound mass.\ The center of a black hole may
    contain a core of maximally compressed bound Aether---a state fundamentally different
    from the mathematical singularity predicted by GR.

  \item The Planck density ($\rho_{\mathrm{Planck}} = c^5 / (\hbar G^2) \approx 5.16 \times 10^{96}\;\mathrm{kg/m^3}$)
    represents the scale at which quantum gravitational effects are expected
    to dominate.\ In the Aether model, this may correspond to the density at which $\Ka$
    nonlinearity becomes significant, naturally limiting the compression and avoiding
    the singularity.
\end{itemize}

Whether the interior reaches a finite equilibrium or undergoes continued collapse depends
on the equation of state of Aether at extreme compression---a regime far beyond current
experimental access but potentially constrained by consistency requirements of the theory.


% ============================================================


% ============================================================
\section{Neutron Stars as Aether Laboratories}
\label{sec:neutron-stars}
% ============================================================

Neutron stars---collapsed stellar cores with masses of $\sim$1.4--2.1\,$M_\odot$
compressed into radii of $\sim$10--13~km---provide the most extreme observable
test of the Aether framework.\ Their central densities ($\sim 10^{17}$--$10^{18}$~kg/m$^3$)
exceed normal nuclear density by factors of 2--8, approaching the compression
regime of the strong-force bridges described in Section~\ref{sec:strong-constraints}.\
Three classes of neutron star observations have direct, quantitative connections
to the theory's predictions.


\subsection{Pulsar Glitches and Quantized Vortex Dynamics}

Pulsars (rotating neutron stars) emit beams of radiation at precisely measured
frequencies.\ Occasionally, a pulsar's rotation frequency increases abruptly by
a fractional amount $\Delta\Omega/\Omega \sim 10^{-9}$--$10^{-6}$---a ``glitch''---then
relaxes back over timescales of hours to months.\ The Vela pulsar glitches
approximately every three years; the 2016 event had a rise time under 12.6~seconds
\cite{palfreyman2018,ashton2019}.

The standard explanation, confirmed by detailed microphysical modeling
\cite{antonelli2025}, involves quantized vortex lines in a superfluid neutron
component:

\begin{enumerate}[leftmargin=2em]
  \item As the crust slows (electromagnetic braking), the interior superfluid
    retains its angular momentum through pinned quantized vortex lines.
  \item Differential rotation builds stress between the pinned vortex array and
    the decelerating crust.
  \item When the stress exceeds a critical threshold, vortex lines unpin in an
    avalanche, transferring angular momentum to the crust---a sudden spin-up.
  \item The system relaxes as vortex-mediated mutual friction recouples the
    components.
\end{enumerate}

In the Aether framework, this is not merely an analogy---it is a macroscopic
manifestation of the same vortex physics that underlies the theory.\ The neutron
star interior is a region where the Aether is compressed to $\sim 10^6$ times
its ambient density, and the superfluid behavior attributed to neutron Cooper
pairs in the standard picture is, in this framework, the superfluid behavior
of the compressed Aether itself.\ The quantized circulation
$\kappa = n\,h/m_A$ that stabilizes electron vortices at ambient density
operates at neutron star densities with the same topological protection---the
vortex core size $\xi$ shrinks (since $c_{\mathrm{local}} = \sqrt{\Ka/\rho}$
decreases with increasing $\rho$), but the quantization condition remains exact.

The glitch phenomenology provides specific observational constraints:

\begin{itemize}[leftmargin=2em]
  \item \textbf{Mutual friction coefficients.}\ The glitch rise time
    ($< 12.6$~s for Vela) and the multi-timescale relaxation constrain the
    mutual friction parameter $\mathcal{B}$, which in the Aether model
    depends on the interaction between quantized vortex cores and the
    surrounding medium at extreme density.

  \item \textbf{Vortex pinning energies.}\ The glitch amplitude and recurrence
    interval constrain the pinning energy per vortex, which in the Aether
    framework is the energy required to detach a quantized circulation
    line from the local density structure of the crustal lattice.

  \item \textbf{Superfluid fraction.}\ Bayesian analysis of the 2016 Vela
    glitch indicates that approximately 70\% of the interior neutrons
    participate in the superfluid component \cite{antonelli2025}.\ In Aether
    terms, this constrains the fraction of the compressed medium that
    maintains quantum coherence and quantized circulation at these extreme
    densities.
\end{itemize}


\subsection{The Equation of State at Extreme Density}

The mass--radius relationship of neutron stars is the most direct observational
probe of the equation of state (EOS) of matter at supranuclear density.\ The
NICER X-ray telescope has measured the radii of several pulsars by tracking
X-ray hotspots on their rotating surfaces, constraining the radius of a
1.4\,$M_\odot$ neutron star to approximately 12.3~km \cite{miller2021,riley2021}.

The Aether framework predicts a specific EOS---the logarithmic equation of state:
%
\begin{equation}
  P = \Ka \ln\!\left(\frac{\rho}{\rho_0}\right) + P_0
\end{equation}
%
derived from the requirement of constant bulk modulus $\Ka$.\ At neutron star
densities ($\rho \sim 10^{17}$--$10^{18}$~kg/m$^3$), the compression ratio
relative to ambient Aether is $\rho/\rhoa \sim 10^6$--$10^7$, approaching but
not reaching the bridge compression ratio $X \approx 7 \times 10^7$.\ The
logarithmic EOS predicts:
%
\begin{equation}
  P_{\mathrm{NS}} \approx \Ka \ln\!\left(\frac{10^{17}}{10^{11}}\right)
    \approx 14\,\Ka
\end{equation}
%
With $\Ka \sim 10^{28}$--$10^{29}$~J/m$^3$, the longitudinal contribution
alone gives $P_{\mathrm{long}} \sim 10^{29}$--$10^{30}$~J/m$^3$---below the
observed neutron star core pressures of $\sim$$10^{34}$--$10^{35}$~Pa by
$\sim$5 orders of magnitude.\ As discussed below, the spin stiffness
$K_{\mathrm{spin}} \propto \rho$ contributes an additional term that
dominates at nuclear density, raising the effective pressure to the correct
range.

The logarithmic EOS is qualitatively different from the polytropic and piecewise
models typically used in nuclear astrophysics.\ Its constant bulk modulus
produces a pressure that increases less steeply than a stiff polytrope
($P \propto \rho^\gamma$ with $\gamma > 2$) at moderate density, but never
softens at high density either---the logarithm continues to rise without limit,
though slowly.\ Whether this specific functional form is consistent with the
observed mass--radius constraints from NICER and the tidal deformability
$\tilde{\Lambda}$ measured from the binary neutron star merger GW170817
\cite{ligo2017gw170817} is a quantitative test that could confirm or falsify
the predicted EOS.

An important caveat: the logarithmic EOS with constant $\Ka$ has been confirmed
experimentally only at compression ratios up to $n^2 \approx 6$ (via the Fresnel
drag coefficient in diamond).\ Neutron star interiors probe compression ratios
of $10^6$--$10^7$, where the EOS has not been independently verified.\
Furthermore, the bridge analysis (Paper~5) identifies a second elastic modulus:
the spin stiffness $K_{\mathrm{spin}} \propto \rho$ governs transverse modes
and grows with density, while the bulk modulus $\Ka$ (constant) governs
longitudinal modes.\ At high compression $X$, the effective stiffness
$K_{\mathrm{eff}} \approx \Ka + \alpha\,K_{\mathrm{spin}} = \Ka(1 + \alpha X)$
is dominated by the spin term, giving an effective speed of sound
$c_{\mathrm{eff}} \approx \sqrt{\alpha}\,c$ at nuclear density
(where $\alpha$ is the spin--compression coupling strength).\ For
$\alpha \approx 0.25$, $c_{\mathrm{eff}} \approx 0.5\,c$---consistent with
neutron star constraints.\ The tidal deformability $\tilde{\Lambda}$ from
GW170817 therefore tests the \emph{combined} bulk and spin stiffness, not
$\Ka$ alone.


\subsection{Deconfinement and Quark Cores}

The central density of the most massive neutron stars ($\gtrsim 2\,M_\odot$)
may exceed the QCD deconfinement threshold of $\sim$5--10 times nuclear density.\
If so, the innermost core would contain deconfined quark matter---a quark--gluon
plasma (QGP).\ Current multi-messenger observations (gravitational waves from
mergers, NICER mass--radius constraints, and nuclear theory benchmarks) are
actively probing whether such phase transitions occur inside neutron stars
\cite{li2025}.

The Aether framework makes a specific prediction for what deconfinement means
physically: the binding bridges (compressed Aether strands anchored by quark
vortex endpoints) dissolve when thermal energy exceeds the elastic energy
per bridge cross-section.\ In a neutron star core, this occurs when the
temperature and density reach the deconfinement regime described in
Section~\ref{sec:strong-force}.

The bridge clarification developed in
Section~\ref{sec:strong-constraints}---that bridges are elastic compressions
of the same superfluid phase, stabilized by quark vortex endpoints---leads to
a distinctive prediction.\ When deconfinement occurs, the elastic compression
energy stored in the bridges ($\varepsilon \approx \Ka\, X$ per unit volume of
bridge material) is released as the compressed Aether relaxes to local
equilibrium.\ This energy release should produce an observable signature:
a softening of the effective EOS at the deconfinement transition, potentially
detectable as a kink in the mass--radius curve or a characteristic frequency
in the gravitational wave spectrum of a neutron star merger.

The theory also predicts that the quark--hadron transition is \emph{not} a
sharp first-order phase transition (as some QCD models suggest) but a gradual
crossover.\ In the Aether picture, deconfinement is the thermal disruption
of individual bridge endpoints---a process that occurs progressively as
temperature and density increase.\ Different bridges dissolve at different
local conditions, producing a smooth transition rather than a discontinuous
jump.\ Recent Bayesian analyses of neutron star data favor smooth crossover
models over sharp phase transitions \cite{li2025}, consistent with this
prediction.



% ============================================================
\section{Summary of Theoretical Properties}
% ============================================================

\subsection{Properties of Black Holes}

\begin{itemize}[leftmargin=2em]
  \item The event horizon is a sonic horizon: the Aether inflow velocity $\vin(r) = \sqrt{2GM/r}$
    reaches $c$ at the Schwarzschild radius $\rs = 2GM/c^2$, preventing outward wave
    propagation.\ This is physically identical to a sonic horizon in fluid dynamics.
  \item Hawking radiation arises from asymmetric transient binding near the horizon: the
    restoring pressure pulse from pair annihilation is partially swept inward by the
    near-$c$ inflow, producing a net outward energy flux.\ The temperature depends on the
    inflow velocity gradient at the horizon: $|d\vin/dr| = c^3/(4GM)$, giving $T \propto 1/M$.
  \item The Hawking effect and the Unruh effect are the same phenomenon in the Aether model:
    asymmetric transient binding caused by velocity gradients in the medium, whether from
    gravitational inflow or from acceleration.
  \item The photon sphere at $r = 3GM/c^2$ arises from the acoustic metric of the inflow
    bending reaching the strength required for circular photon orbits.\ The black hole
    shadow observed by the Event Horizon Telescope is the photon sphere's projection.
  \item Frame dragging (Lense--Thirring effect) is the co-rotation of Aether driven by the
    aggregate electron vortex circulation of a rotating mass.\ A rotating black hole
    creates an ergosphere where the dragged Aether flow exceeds what any stationary object
    can resist.
  \item The no-hair theorem follows from the sonic horizon: only mass (inflow profile),
    angular momentum (rotational flow), and charge (radial flow direction) produce long-range Aether patterns that survive outside the horizon.\ All other information is swept
    inward.
  \item The singularity may be avoided if $\Ka$ increases at extreme compression, limiting the
    maximum achievable density to a finite value (potentially at or near the Planck
    density).
\end{itemize}




\subsection{Properties of Neutron Stars}

\begin{itemize}[leftmargin=2em]
  \item Neutron star interiors are regions where Aether is compressed to
    $\sim 10^6$--$10^7$ times ambient density, approaching the bridge
    compression ratio $X \approx 7 \times 10^7$.
  \item Pulsar glitches are macroscopic manifestations of quantized vortex
    dynamics: the same circulation quantization $\kappa = nh/m_A$ that
    stabilizes electrons at ambient density operates in the extreme-density
    interior, with vortex unpinning avalanches producing sudden spin-ups.
  \item The logarithmic equation of state $P = \Ka \ln(\rho/\rho_0) + P_0$,
    combined with the spin stiffness $K_{\mathrm{spin}} \propto \rho$ at
    nuclear density, produces neutron star core pressures consistent with
    observations ($\sim 10^{34}$--$10^{35}$~Pa) and makes specific,
    falsifiable predictions for the mass--radius relationship and tidal
    deformability.
  \item At sufficiently high core density and temperature, the Aether bridges
    confining quarks dissolve (deconfinement), producing a quark core.\ The
    theory predicts this transition is a gradual crossover (progressive bridge
    disruption) rather than a sharp first-order phase transition.
  \item The released elastic compression energy from bridge dissolution should
    produce a softening of the effective EOS at the deconfinement transition,
    potentially detectable in gravitational wave spectra from neutron star mergers.
\end{itemize}

\subsection{Properties of Dark Energy and Dark Matter}

\begin{itemize}[leftmargin=2em]
  \item Dark energy is the thermal pressure of hot Aether in cosmic voids,
    pushing against the cold Aether regions near matter.\ Driven by
    accumulated radiated energy from stars with no mass formation to absorb it.\ Produces
    real expansion and real Doppler redshift, naturally explaining supernova time dilation.
  \item Dark matter effects arise from the two-fluid first-sound correction: the
    normal fraction increases the first-sound speed $c_1 > c$, and the outward gradient
    $dc_1^2/dr > 0$ provides additional centripetal acceleration through the acoustic
    metric---no additional mass required.
  \item The Aether thermal pressure model predicts that dark energy evolves
    over cosmic time---consistent with DESI observations of a time-varying
    equation of state at 2.8--4.2$\sigma$ significance.
  \item Photons traversing voids encounter modified Aether properties
    (higher temperature, larger normal fraction), potentially amplifying
    the integrated Sachs--Wolfe effect beyond $\Lambda$CDM predictions.
  \item Void density profiles are predicted to be approximately universal
    when scaled by void radius, due to the invariance of $\Ka$.
\end{itemize}



% ============================================================

% ============================================================
\section{Experimental Support and Connections to Established Physics}
% ============================================================

The following experimentally confirmed phenomena support or align with the aspects of the theory presented in this paper:

\begin{enumerate}[leftmargin=2em]
  \item \textbf{Accelerating expansion.} The observed acceleration of cosmic expansion is consistent
    with this theory's prediction of a positive feedback loop: matter radiates energy
    into voids, heating the Aether, increasing thermal pressure, and driving faster
    expansion.

  \item \textbf{Supernova time dilation.} Distant supernova light curves are stretched proportional to
    their redshift.\ This is naturally produced by the real expansion mechanism of this
    theory (thermal pressure driving genuine recession), matching the predictions of
    standard expansion-based cosmology.

  \item \textbf{Failed dark matter particle searches.} Extensive experimental programs (LUX \cite{lux2017}, XENON \cite{xenon2018},
    PandaX \cite{pandax2017}, and others) have failed to detect dark matter particles despite decades of
    effort.\ This is consistent with this theory's prediction that dark matter is not a
    particle but a medium property---the gravitating kinetic energy of the
    Vinen-regulated fountain counterflow in the two-fluid Aether.

  \item \textbf{Black hole shadow observations.} The Event Horizon Telescope imaged the shadow of
    M87* \cite{eht2019} and Sagittarius~A* \cite{eht2022}, confirming the existence and size of the photon
    sphere to approximately 10\% precision.\ Since the Aether model
    reproduces the weak-field bending angle exactly, the strong-field photon sphere and
    shadow size are expected to be consistent with these observations, subject to
    quantitative verification through full ray-tracing calculations.

  \item \textbf{Acoustic Hawking radiation \cite{steinhauer2019}.} A Bose--Einstein condensate experiment
    created a sonic horizon in a flowing superfluid and observed thermal phonon emission
    at the predicted Hawking temperature.\ In the Aether model, this is not an analogy but
    the same physical mechanism: thermal radiation from asymmetric fluctuations at a sonic
    horizon in a flowing medium.


  \item \textbf{Analog Hawking radiation in a Bose--Einstein condensate.}\
    Steinhauer \cite{steinhauer2016} created a sonic black hole in an ultracold
    rubidium condensate by engineering a flow exceeding the local speed of sound,
    and observed spontaneous emission of entangled phonon pairs at the acoustic
    horizon---the analog of Hawking radiation.\ In the Aether model, black holes
    are sonic horizons where Aether inflow velocity reaches $c$.\ The observation
    that a superfluid sonic horizon spontaneously produces quantum radiation
    confirms the mechanism by which this theory predicts Hawking emission.

  \item \textbf{Rotating black hole geometry from a giant quantum vortex.}\
    Svancara et al.\ \cite{svancara2024} stabilized a giant quantum vortex in
    superfluid ${}^4$He carrying thousands of circulation quanta and showed that
    surface wave dynamics reproduce the curved spacetime geometry of a rotating
    (Kerr) black hole, including ergosphere signatures and ringdown physics.\
    In the Aether model, gravity \emph{is} Aether flow, and a black hole
    \emph{is} a sonic horizon.\ The fact that a large vortex in a quantum
    superfluid spontaneously generates the geometry of a rotating black hole
    is a direct prediction of this framework.

  \item \textbf{DESI evidence for evolving dark energy.}\ The Dark Energy
    Spectroscopic Instrument \cite{desi2025}, using baryon acoustic oscillation
    measurements from over 14 million extragalactic objects, finds 2.8--4.2$\sigma$
    evidence that the dark energy equation of state varies with redshift---
    inconsistent with a cosmological constant and consistent with the Aether
    model's prediction that dark energy (thermal void pressure) evolves over
    cosmic time.

  \item \textbf{Anomalous ISW signal from supervoids.}\ Galaxy surveys find
    that supervoids produce integrated Sachs--Wolfe temperature imprints
    approximately $5.2 \pm 1.6$ times stronger than $\Lambda$CDM predicts
    \cite{kovacs2022}.\ The Aether model's prediction of modified photon
    propagation through high-temperature void regions (larger normal fraction) could
    account for this amplification.

  \item \textbf{Universal void density profiles.}\ When scaled by the void
    effective radius, observed void density profiles are approximately
    universal---independent of void size and cosmic epoch---consistent with
    a medium whose bulk modulus $\Ka$ is invariant everywhere.


  \item \textbf{Pulsar glitches as quantized vortex avalanches.}\ The Vela
    pulsar's 2016 glitch \cite{palfreyman2018,ashton2019} exhibited a rise
    time under 12.6~seconds and multi-timescale relaxation, both explained
    by the sudden unpinning of quantized vortex lines in a superfluid
    interior.\ In the Aether model, these are quantized circulation lines
    ($\kappa = nh/m_A$) in extreme-density Aether---the same vortex physics
    that produces the Pauli exclusion principle and electron stability,
    observed at macroscopic scales.

  \item \textbf{NICER mass--radius constraints.}\ X-ray observations of
    millisecond pulsars \cite{miller2021,riley2021} constrain neutron star
    radii to $\sim$12.3~km for 1.4\,$M_\odot$ stars, directly testing the
    equation of state at supranuclear densities.\ The logarithmic EOS
    predicted by the Aether framework ($P = \Ka \ln(\rho/\rho_0) + P_0$)
    combined with the spin stiffness $K_{\mathrm{spin}} \propto \rho$,
    produces neutron star core pressures consistent with observations
    and can be quantitatively compared to these constraints.

  \item \textbf{Binary neutron star merger GW170817.}\ The tidal
    deformability measured from this merger \cite{ligo2017gw170817}
    constrains the stiffness of the EOS at 2--6 times nuclear density.
    The Aether prediction of constant $\Ka$ (logarithmic EOS) makes a
    specific, falsifiable prediction for the tidal deformability parameter
    $\tilde{\Lambda}$.

  \item \textbf{Galaxy rotation curves.}\ Galaxies exhibit flat rotation
    curves requiring $\rho_{\mathrm{DM}} \propto 1/r^2$ halos with cored
    central profiles.\ The two-fluid first-sound correction
    (Section~\ref{sec:dark-matter}) naturally produces isothermal
    $1/r^2$ profiles with cores, matching observations without invoking
    a dark matter particle.

  \item \textbf{The Hubble tension.}\ Early-universe (CMB) and
    late-universe (distance ladder) measurements of $H_0$ disagree at
    $\sim$5$\sigma$.\ The Aether framework resolves this through
    environment-dependent thermal pressure: voids, which dominate the
    local distance ladder path, have higher Aether temperature and
    therefore faster local expansion than the cosmic mean sampled by
    the CMB.

\end{enumerate}


% ============================================================
\section{Open Questions and Testable Predictions}
% ============================================================

The following areas relevant to this paper remain underdeveloped and represent opportunities for further refinement or falsification:

\begin{enumerate}[leftmargin=2em]
  \item \textbf{Void temperature and expansion rate.} The theory predicts a direct relationship between
    void temperature (or thermal energy content) and expansion rate.\ Measurements
    correlating the thermal properties of cosmic voids with local expansion rates could
    test whether thermal pressure is the driver of dark energy.

  \item \textbf{Two-fluid first-sound dark matter: open calculations.}\
    The first-sound correction mechanism (Section~\ref{sec:dark-matter},
    Eq.~\ref{eq:first-sound}) has the correct sign and produces naturally
    cored profiles, but its magnitude depends on the Aether's excitation
    spectrum.\ The critical open calculations are:
    (a)~determine the Aether's excitation spectrum beyond the phonon
    regime---does it have roton-like intermediate-wavelength modes that
    break the temperature-independence of the phonon gas correction?\
    For a pure phonon gas, the first-sound correction is $c^2/3$
    (constant, producing no spatial gradient); non-phonon excitations
    are needed to create the temperature-dependent variation required
    for dark matter ($\Delta c_1^2/c^2 \sim 10^{-6}$ for galaxies;
    the ideal BEC gives $\sim$$10^{-4}$, $\sim$$50\times$ too large);
    (b)~derive the three-dimensional temperature profile $T(r,z)$ from
    first principles (stellar heating vs.\ binding cooling vs.\ conduction);
    (c)~determine whether the resulting $c_1(r)$ profile produces flat
    rotation curves;
    (d)~calculate the CMB power spectrum from the cosmological superfluid
    phase transition;
    (e)~verify Bullet Cluster consistency.

  \item \textbf{Quantitative Aether-enhanced ISW effect.}\ The Aether
    framework predicts that photons traversing cosmic voids encounter
    modified medium properties (higher temperature, larger normal fraction).
    Can a quantitative calculation of the resulting temperature
    imprint reproduce the factor-of-five amplification of the ISW signal
    observed in supervoids \cite{kovacs2022}?\ This requires modeling
    the void temperature and density profiles within the three-property
    framework.

  \item \textbf{Dark energy equation of state from the two-fluid model.}\
    The normal-component bulk viscosity $\zeta_n \sim 3 \times 10^7$~Pa\,s
    (Eq.~\ref{eq:zeta-normal}) and the cosmic star formation rate history
    together determine the heating--dilution balance that governs the dark
    energy equation of state $w(z)$.\ Can the predicted $w_0$--$w_a$
    match the DESI observation of a phantom-to-quintessence crossing at
    $z \approx 0.5$?\ This is a specific numerical calculation using
    $\dot\varepsilon_{\mathrm{heat}} = L_*(z) +
    9\,\zeta_n\,H^2(z)$ and $\dot\varepsilon_{\mathrm{dilute}}
    = 3H\,P_{\mathrm{th}}$.

  \item \textbf{Cosmic coincidence from shared temperature structure.}\
    Both dark matter (the first-sound correction $\Delta c_1^2/c^2$)
    and dark energy (the thermal pressure $P_{\mathrm{th}}$) arise from
    the same temperature field $T(\mathbf{x})$.\ Does the observed ratio
    $\Omega_{\mathrm{DM}}/\Omega_\Lambda \approx 0.4$ follow naturally
    from the temperature profile shape, without fine-tuning?\ If so, this
    would resolve the cosmic coincidence problem.

  \item \textbf{Environment-dependent expansion and dark matter
    correlation.}\ The two-fluid framework predicts that the local
    expansion rate and the local dark matter fraction are correlated
    (both increase with the normal fraction, which increases with
    temperature).\ Upcoming DESI and Euclid void catalogs could test
    whether galaxies in hotter environments (near large voids) show both
    larger ``dark matter'' halos and faster local expansion.

  \item \textbf{Bulk viscosity $\zeta_n$ and observational
    constraints.}\ The predicted $\zeta_n \sim 10^6$--$3 \times 10^7$~Pa\,s
    gives a kinematic bulk viscosity
    $\nu_\zeta = \zeta_n/\rho_n \sim 10^{-4}$--$10^{-3}$~m$^2$/s
    (tiny, because $\rho_n \sim 10^{10}$~kg/m$^3$).\ The viscous damping
    scale is therefore submicroscopic---the normal fraction dissipates
    energy from cosmic expansion without damping astrophysical density
    perturbations.\ In superfluids, $\zeta$ arises from three-phonon
    relaxation and diverges near $T_c$; can $\zeta_n$ be independently
    constrained from the Aether's excitation spectrum or from its effect
    on second-sound attenuation?

  \item \textbf{Void temperature ratio and the Hubble tension.}\
    The quantitative analysis of Section~1.7 shows that matching the
    observed Hubble tension requires a void thermal pressure excess of
    $\eta \approx 0.22$ (Eq.~\ref{eq:eta-required}), corresponding to a
    void-to-wall temperature ratio $T_{\mathrm{void}}/T_{\mathrm{wall}}
    \approx 2.5$ (Eq.~\ref{eq:T-ratio}).\ Can this ratio be derived from a
    thermodynamic model of void heating and wall cooling over cosmological
    time?\ Specifically: (a)~does the integrated radiative energy deposited
    into voids, combined with energy absorbed by mass formation in walls,
    produce the required factor of $\sim$2.5?\ (b)~Does the predicted
    $H_0$ anisotropy pattern match the dipole and quadrupole structure seen
    in recent anisotropy analyses of the distance ladder?\ (c)~Can the void
    abundance predicted by the thermal pressure mechanism be tested against
    upcoming DESI and Euclid void catalogs?

  \item \textbf{Two-fluid Aether: critical temperature and superfluid
    fraction.}\ The two-fluid dark matter mechanism requires a finite
    $T_c$ and a continuously varying superfluid fraction.\ For
    $m_A \approx 500$--$850$~MeV/$c^2$ (Paper~7), the ideal BEC critical
    temperature is $T_c \sim 10^{10}$--$10^{11}$~K, placing the cosmological
    phase transition well before nucleosynthesis.\ Key questions: (a)~What is
    the superfluid fraction in cosmic voids today?\ (b)~Does the
    anomalous specific heat near $T_c$ (the lambda anomaly)
    thermodynamically stabilize the halo boundary, creating sharper
    edges than NFW profiles?\ (c)~Is $T_c$ high enough that the
    entire present-day universe remains below it ($T_{\mathrm{void}}
    < T_c$), ensuring the condensate exists everywhere and the universe
    is transparent?

  \item \textbf{Photon dispersion from the normal fraction in voids.}\
    At macroscopic wavelengths, photons propagate at the two-fluid
    first-sound speed $c_1$ (via spin-orbit coupling), which is
    achromatic.\ However, at wavelengths approaching the spin healing
    length $\xi_s$, the spin-wave dispersion relation acquires
    corrections.\ Because the normal fraction is larger in voids, the
    spin-wave--normal-fluid coupling (suppressed by
    $P_{\mathrm{th}}/\Ka \sim 10^{-38}$) should produce a residual
    \emph{differential} dispersion: photons traversing supervoids
    should exhibit slightly more dispersion than photons traveling
    through matter-rich regions (smaller normal fraction).\ The Fermi
    LAT constrains Lorentz-violating dispersion to energy scales above
    $7.6\,E_{\mathrm{Planck}}$ \cite{vasileiou2013}, but a differential
    void-versus-filament comparison could reveal or constrain the void
    superfluid fraction.

  \item \textbf{Matter formation at the cosmological phase transition.}\
    During the Aether's superfluid phase transition ($T_c \sim
    10^9$~K, $t \sim 1$~s), the Kibble--Zurek mechanism
    \cite{bauerle1996,ruutu1996} predicts copious creation of
    quantized vortices (= particles) at the transition front.\ The
    resulting matter concentrations produce gravitational wells that
    seed structure formation \emph{before} any galaxies exist,
    potentially resolving the early-universe dark matter requirement.\
    Does the spatial spectrum of Kibble--Zurek defects from the Aether
    transition match the CMB power spectrum?

  \item \textbf{Second sound in the cosmic Aether.}\ The two-fluid
    model predicts a second propagation mode: temperature waves
    (second sound) at speed $c_2 \ll c$.\ Sudden heating events
    (supernovae, galaxy mergers, quasar ignition) would launch
    second-sound pulses through the Aether.\ These temperature waves
    would appear as propagating oscillations in the local dark energy
    density.\ Could second sound contribute to the baryon acoustic
    oscillation signal or other large-scale features?\ The speed
    $c_2$ is determined by the superfluid fraction and the Aether's
    thermodynamic properties---both calculable in principle.

  \item \textbf{Critical velocity and quantum turbulence near compact
    objects.}\ The gravitational inflow $\vin = \sqrt{2GM/r}$
    exceeds the Landau critical velocity $v_c = \Delta/p_{\mathrm{roton}}$
    at some radius near compact objects, nucleating quantized vortices
    and driving the Aether into a quantum turbulent state.\ For
    neutron stars, the resulting vortex dynamics may connect to the
    observed pulsar glitch mechanism (vortex unpinning events).\ For
    black holes, the region between $v_c$ and the sonic horizon is a
    quantum turbulent shell whose properties affect the Hawking
    temperature and gravitational wave ringdown signatures.

\end{enumerate}



% ============================================================
\section{Proposed Experiments}
% ============================================================

The following experiments target specific predictions relevant to this paper:

\begin{enumerate}[leftmargin=2em]
  \item \textbf{Neutron star gravitational redshift as a strong-field test.}\
    General relativity predicts the surface gravitational redshift of a
    neutron star as $z = (1 - 2GM/Rc^2)^{-1/2} - 1$.\ The Aether model's
    logarithmic equation of state produces a nonlinear density profile in the
    strong-field regime ($\delta\rho/\rho_0 \gtrsim 0.1$) that may yield a
    different redshift for the same mass and radius.\ NICER pulse-profile
    modeling now provides simultaneous $M$ and $R$ measurements for several
    millisecond pulsars; combining these with spectroscopic redshift
    measurements from X-ray absorption lines (e.g., 1E~1207.4$-$5209,
    EXO~0748$-$676) would allow a direct comparison between the GR and Aether
    predictions.\ A statistically significant discrepancy at the
    ${\sim}\,5$--$10\%$ level would distinguish the two frameworks in the
    strong-field regime.

  \item \textbf{CMB frame and Aether rest frame coincidence.}\
    Special relativity asserts that no preferred rest frame exists; the CMB
    dipole (our motion at ${\sim}\,370$~km/s relative to the cosmic microwave
    background) is merely a Doppler effect with no deeper significance.\ The
    Aether model predicts that the medium \emph{does} have a rest frame.\ If
    the Aether rest frame coincides with the CMB rest frame (natural if CMB
    photons are Aether excitations), then all precision tests of Lorentz
    invariance should show residual anisotropies aligned with the CMB dipole
    direction.\ If they do \emph{not} coincide, a ``dark flow''---a
    systematic velocity of the Aether relative to the CMB---would produce
    unexplained large-scale anisotropies in galaxy surveys.\ Current
    large-scale structure surveys (DESI, Euclid) and next-generation CMB
    experiments (CMB-S4) have the sensitivity to detect or rule out bulk
    flows at the ${\sim}\,100$~km/s level, constraining or confirming the
    Aether frame hypothesis.

\end{enumerate}

\noindent\textit{For all proposed experiments, see Paper~1: General Overview \cite{otto-overview}.}


% ============================================================
\section{Mathematical Summary}
% ============================================================

This section collects the key equations relevant to the topics covered in this paper.
\noindent\textit{For the complete mathematical summary, see Paper~1: General Overview \cite{otto-overview}.}

\subsection{Fundamental Aether Parameters}
% --------------------------------------------------

The theory posits four microscopic parameters from which all macroscopic
phenomena derive:
%
\begin{center}
\begin{tabular}{@{}llll@{}}
\hline
\textbf{Symbol} & \textbf{Name} & \textbf{Constrained range} & \textbf{Primary constraint} \\
\hline
$\rhoa$  & Ambient density     & $10^{11}$--$10^{12}$~kg/m$^3$         & Electron vortex energy \\
$\Ka$    & Bulk modulus         & $10^{28}$--$10^{29}$~J/m$^3$          & $\Ka = \rhoa\,c^2$ \\
$m_A$    & Aether quantum mass  & $500$--$850$~MeV/$c^2$                & Flux tube width / glueball mass \\
$\xi$    & Healing length       & $0.17$--$0.28$~fm                     & Bridge diameter (lattice QCD) \\
\hline
\end{tabular}
\end{center}

These four are not independent.\ Two defining relations connect them:
%
\begin{align}
  c &= \sqrt{\Ka / \rhoa}
    \label{eq:sum-speed-of-light} \\
  \xi &= \frac{\hbar}{m_A\,c\,\sqrt{2}}
    \label{eq:sum-healing-length}
\end{align}
%
leaving two free Aether parameters (e.g., $\rhoa$ and $m_A$) from which all
others follow.\ (A third parameter, $\epsilon$, enters through gravity; see below.)


% --------------------------------------------------
\subsection{Equation of State}
% --------------------------------------------------

The Aether obeys a logarithmic equation of state, the unique form that produces
a density-independent bulk modulus (Section~\ref{sec:qss}):
%
\begin{equation}
  P = \Ka \ln\!\Bigl(\frac{\rho}{\rho_0}\Bigr) + P_0
  \label{eq:sum-eos}
\end{equation}
%
Key consequence: the speed of sound $c_s = \sqrt{\Ka/\rho}$ varies with
density.\ At ambient density, $c_s = c$.\ Note: in gravitational fields,
Bernoulli's equation for free-fall inflow gives $\Delta\rho/\rhoa = 0$
exactly, so $c$ does not vary near gravitating masses; the density
dependence is relevant for material media and for the two-fluid
cosmological structure.\ This equation of state guarantees that $\Ka$ is
constant under compression, which is experimentally confirmed by:
%
\begin{itemize}[leftmargin=2em]
  \item Velocity time dilation measurements spanning $10^{-6}\,c$ to $0.9994\,c$
    (precision $\sim 10^{-8}$).
  \item The Fresnel drag coefficient $1 - 1/n^2$, exact across materials with
    density ratios up to $\sim$6:1.
\end{itemize}


% --------------------------------------------------
\subsection{Cosmology}
% --------------------------------------------------

\paragraph{Dark energy.}
Thermal pressure of hot Aether in cosmic voids (density essentially uniform,
$\delta\rhoa/\rhoa \sim 10^{-4}$; temperature is the cosmologically relevant variable):
%
\begin{equation}
  P_{\mathrm{void}} = \Ka\,\ln\!\Bigl(\frac{\rho_{\mathrm{void}}}{\rho_0}\Bigr)
    + P_{\mathrm{th}}(T_{\mathrm{void}})
  \label{eq:sum-dark-energy}
\end{equation}
%
where $P_{\mathrm{th}}(T)$ is the thermal pressure from the Aether's
excitation spectrum ($P_{\mathrm{th}} \propto T$ for an ideal gas,
$P_{\mathrm{th}} = \frac{1}{3}a_R T^4$ for a radiation-dominated spectrum;
the actual form is determined by the same excitation spectrum that sets
the dark matter magnitude and the bulk viscosity).\ The theory predicts dark
energy evolves over cosmic time (consistent with DESI observations at
2.8--4.2$\sigma$).

\paragraph{Dark matter.}
The Landau two-fluid first-sound correction: the normal fraction increases
$c_1$ above $c$, and the outward gradient $dc_1^2/dr > 0$ provides additional
centripetal acceleration through the acoustic metric (no additional mass
required).\ Rotation curves and gravitational lensing are automatically
consistent (same acoustic metric).\ The Milky Way halo requires
$\Delta c_1^2/c^2 \approx 6 \times 10^{-6}$; the magnitude depends on
the Aether's excitation spectrum.\ Naturally produces cored profiles.

\paragraph{Vacuum catastrophe resolution.}
The QFT vacuum energy density $\sim 10^{113}$~J/m$^3$ is reinterpreted as
$\Ka \sim 10^{28}$~J/m$^3$ (the actual ground-state energy density of the
Aether), resolving the $10^{85}$-fold discrepancy between the predicted vacuum energy and the Aether ground-state energy density (compared to the $10^{120}$-fold discrepancy with the observed cosmological constant).


% --------------------------------------------------

\section*{Acknowledgments}
\addcontentsline{toc}{section}{Acknowledgments}

The author acknowledges the use of Claude (Anthropic) for \LaTeX{} formatting
and editorial assistance.

\bigskip


% ============================================================
% BIBLIOGRAPHY
% ============================================================

\begin{thebibliography}{99}

\bibitem{hawking1974}
Hawking, S.~W.,
``Black hole explosions?''
\textit{Nature}, \textbf{248}, 30--31 (1974).


\bibitem{hawking1975}
Hawking, S.~W.,
``Particle creation by black holes,''
\textit{Communications in Mathematical Physics}, \textbf{43}, 199--220 (1975).


\bibitem{steinhauer2019}
Mu\~{n}oz de Nova, J.~R., Golubkov, K., Kolobov, V.~I., and Steinhauer, J.,
``Observation of thermal Hawking radiation and its temperature in an analogue black hole,''
\textit{Nature}, \textbf{569}, 688--691 (2019).


\bibitem{everitt2011}
Everitt, C.~W.~F., DeBra, D.~B., Parkinson, B.~W., et al.,
``Gravity Probe B: Final Results of a Space Experiment to Test General Relativity,''
\textit{Physical Review Letters}, \textbf{106}, 221101 (2011).


\bibitem{eht2019}
Event Horizon Telescope Collaboration,
``First M87 Event Horizon Telescope Results.\ I.\ The Shadow of the Supermassive Black Hole,''
\textit{The Astrophysical Journal Letters}, \textbf{875}, L1 (2019).


\bibitem{eht2022}
Event Horizon Telescope Collaboration,
``First Sagittarius~A* Event Horizon Telescope Results.\ I.\ The Shadow of the Supermassive Black Hole at the Center of the Milky Way,''
\textit{The Astrophysical Journal Letters}, \textbf{930}, L12 (2022).


\bibitem{lux2017}
Akerib, D.~S.\ et al.\ (LUX Collaboration),
``Results from a search for dark matter in the complete LUX exposure,''
\textit{Physical Review Letters}, \textbf{118}, 021303 (2017).


\bibitem{xenon2018}
Aprile, E.\ et al.\ (XENON Collaboration),
``Dark Matter Search Results from a One Ton-Year Exposure of XENON1T,''
\textit{Physical Review Letters}, \textbf{121}, 111302 (2018).


\bibitem{pandax2017}
Cui, X.\ et al.\ (PandaX-II Collaboration),
``Dark Matter Results From 54-Ton-Day Exposure of PandaX-II Experiment,''
\textit{Physical Review Letters}, \textbf{119}, 181302 (2017).


\bibitem{penrose1969}
Penrose, R.,
``Gravitational Collapse: The Role of General Relativity,''
\textit{Rivista del Nuovo Cimento}, \textbf{1}, 252--276 (1969).


\bibitem{lense1918}
Lense, J.\ and Thirring, H.,
``\"{U}ber den Einfluss der Eigenrotation der Zentralk\"{o}rper auf die Bewegung der Planeten und Monde nach der Einsteinschen Gravitationstheorie,''
\textit{Physikalische Zeitschrift}, \textbf{19}, 156--163 (1918).


\bibitem{hubble1929}
Hubble, E.,
``A relation between distance and radial velocity among extra-galactic nebulae,''
\textit{Proc.\ National Academy of Sciences}, \textbf{15}, 168--173 (1929).


\bibitem{vasileiou2013}
Vasileiou, V., Jacholkowska, A., Piron, F., et al.,
``Constraints on Lorentz invariance violation from Fermi--Large Area Telescope
observations of gamma-ray bursts,''
\textit{Phys.\ Rev.\ D}, \textbf{87}, 122001 (2013).


\bibitem{desi2025}
DESI Collaboration,
``DESI DR2 Results II: Measurements of baryon acoustic oscillations and cosmological constraints,''
\textit{Phys.\ Rev.\ D}, \textbf{112}, 083515 (2025).


\bibitem{des2026}
Dark Energy Survey Collaboration,
``Dark Energy Survey Year 6 Results: Cosmological Constraints from Galaxy Clustering and Weak Lensing,''
arXiv:2601.14559 (2026).


\bibitem{hohmann2026}
Hohmann, M., Pfeifer, C., and Voicu, N.,
``Cosmological Finsler spacetimes and observational constraints,''
\textit{J.\ Cosmol.\ Astropart.\ Phys.}, 2026.


\bibitem{wiltshire2025}
Seifert, A., Lane, Z.~G., Galoppo, M., Ridden-Harper, R., and Wiltshire, D.~L.,
``Supernovae evidence for foundational change to cosmological models,''
\textit{Mon.\ Not.\ R.\ Astron.\ Soc.\ Lett.}, \textbf{537}, L55--L60 (2025).


\bibitem{vielva2004}
Vielva, P., et al.,
``Detection of non-Gaussianity in the WMAP first-year data using spherical wavelets,''
\textit{Astrophys.\ J.}, \textbf{609}, 22--34 (2004).


\bibitem{szapudi2015}
Szapudi, I., et al.,
``Detection of a supervoid aligned with the cold spot of the cosmic microwave background,''
\textit{Mon.\ Not.\ R.\ Astron.\ Soc.}, \textbf{450}, 288--294 (2015).


\bibitem{kovacs2022}
Kov\'acs, A., et al.,
``The DES view of the Eridanus supervoid and the CMB cold spot,''
\textit{Mon.\ Not.\ R.\ Astron.\ Soc.}, \textbf{510}, 216--229 (2022).



\bibitem{steinhauer2016}
Steinhauer, J.,
``Observation of quantum Hawking radiation and its entanglement in an analogue black hole,''
\textit{Nature Phys.}, \textbf{12}, 959--965 (2016).


\bibitem{svancara2024}
Svancara, P., Smaniotto, P., Solidoro, L., et al.,
``Rotating curved spacetime signatures from a giant quantum vortex,''
\textit{Nature} (2024). DOI: 10.1038/s41586-024-07176-8.


\bibitem{bauerle1996}
B\"auerle, C., Bunkov, Yu.~M., Fisher, S.~N., Godfrin, H., and Pickett, G.~R.,
``Laboratory simulation of cosmic string formation in the early Universe using superfluid $^3$He,''
\textit{Nature}, \textbf{382}, 332--334 (1996).


\bibitem{ruutu1996}
Ruutu, V.~M.~H., Eltsov, V.~B., Gill, A.~J., et al.,
``Vortex formation in neutron-irradiated superfluid $^3$He as an analogue of cosmological defect formation,''
\textit{Nature}, \textbf{382}, 334--336 (1996).


\bibitem{palfreyman2018}
  Palfreyman, J., Dickey, J.~M., Hotan, A., Ellingsen, S., van Straten, W.,
  ``A glitch in the millisecond pulsar J0613-0200,''
  \textit{Nature}, \textbf{556}, 219--222 (2018).


\bibitem{ashton2019}
  Ashton, G., Lasky, P.~D., Graber, V., Palfreyman, J.,
  ``Rotational evolution of the Vela pulsar during the 2016 glitch,''
  \textit{Nature Astronomy}, \textbf{3}, 1143--1148 (2019).


\bibitem{antonelli2025}
  Antonelli, M., Montoli, A., Pizzochero, P.~M.,
  ``A microphysical probe of neutron star interiors:
  constraining the equation of state with glitch dynamics,''
  \textit{arXiv preprint} 2510.20791 (2025).


\bibitem{miller2021}
  Miller, M.~C., et~al.,
  ``The radius of PSR J0740+6620 from NICER and XMM-Newton data,''
  \textit{The Astrophysical Journal Letters}, \textbf{918}, L28 (2021).


\bibitem{riley2021}
  Riley, T.~E., et~al.,
  ``A NICER view of the massive pulsar PSR J0740+6620
  informed by radio timing and XMM-Newton spectroscopy,''
  \textit{The Astrophysical Journal Letters}, \textbf{918}, L27 (2021).


\bibitem{ligo2017gw170817}
  Abbott, B.~P., et~al.\ (LIGO/Virgo Collaboration),
  ``GW170817: observation of gravitational waves from a
  binary neutron star inspiral,''
  \textit{Physical Review Letters}, \textbf{119}, 161101 (2017).


\bibitem{li2025}
  Li, A., et~al.,
  ``Simultaneously constraining the neutron star equation of state
  and mass distribution through multimessenger observations
  and nuclear benchmarks,''
  \textit{Physical Review C}, \textbf{111}, 025806 (2025).



\bibitem{planck2020}
Planck Collaboration,
``Planck 2018 results. VI. Cosmological parameters,''
\textit{Astron.\ Astrophys.}, \textbf{641}, A6 (2020).


\bibitem{riess2022}
Riess, A.~G., et~al.,
``A comprehensive measurement of the local value of the Hubble constant with
1~km\,s$^{-1}$\,Mpc$^{-1}$ uncertainty from the Hubble Space Telescope and
the SH0ES team,''
\textit{Astrophys.\ J.\ Lett.}, \textbf{934}, L7 (2022).


\bibitem{wong2020}
Wong, K.~C., et~al.,
``H0LiCOW XIII.\ A 2.4\% measurement of $H_0$ from lensed quasars:
$5.3\sigma$ tension between early- and late-Universe probes,''
\textit{Mon.\ Not.\ R.\ Astron.\ Soc.}, \textbf{498}, 1420--1439 (2020).


\bibitem{haslbauer2020}
Haslbauer, M., Banik, I., and Kroupa, P.,
``The KBC void and Hubble tension contradict $\Lambda$CDM on a Gpc scale---Milgromian
dynamics as a timely alternative,''
\textit{Mon.\ Not.\ R.\ Astron.\ Soc.}, \textbf{499}, 2845--2883 (2020).


\bibitem{volovik2003}
Volovik, G.~E.,
\textit{The Universe in a Helium Droplet}
(Oxford University Press, Oxford, 2003).

\bibitem{berezhiani2015}
Berezhiani, L. and Khoury, J.,
``Theory of dark matter superfluidity,''
\textit{Phys. Rev. D}, \textbf{92}, 103510 (2015).
[arXiv:1507.01019]

\bibitem{otto-gravity}
Otto, E.,
``Unifying Theory of Dynamic Aether Mechanics: Gravity,''
(2026).

\bibitem{otto-overview}
Otto, E.,
``Unifying Theory of Dynamic Aether Mechanics: General Overview,''
(2026).

\bibitem{otto-quantum}
Otto, E.,
``Unifying Theory of Dynamic Aether Mechanics: Quantum Mechanics,''
(2026).

\bibitem{fixsen1996}
Fixsen, D.~J., Cheng, E.~S., Gales, J.~M., et al.,
``The Cosmic Microwave Background Spectrum from the Full COBE FIRAS Data Set,''
\textit{Astrophys.\ J.}, \textbf{473}, 576--587 (1996).
\end{thebibliography}

\end{document}